\section{Background and Motivation}

The global transition towards electrified transportation, renewable
energy integration, and grid-scale storage has made the lithium-ion
battery one of the most strategically important technologies of the
twenty-first century.  Despite decades of engineering progress,
lithium-ion cells remain intrinsically fragile: their electrochemical
stability depends on maintaining terminal voltage, current, and
temperature within narrow operating windows.  Violations of these
limits---even momentary---accelerate solid-electrolyte interphase
(SEI) growth, lithium plating, mechanical particle cracking, and, in
severe cases, thermal runaway \cite{vetter2005ageing}.

A \emph{Battery Management System} (BMS) is the firmware and hardware
subsystem responsible for enforcing these limits in real time while
simultaneously optimising competing objectives: charging as fast as
possible (to minimise user inconvenience), maintaining cell
temperature within acceptable bounds (to extend cycle life), limiting
degradation-inducing high-current and high-SOC exposure (to maximise
capacity retention), and guaranteeing hard electrochemical safety
constraints (to prevent cell failure or fire).

These four objectives are fundamentally in tension.  Faster charging
requires higher current, which generates more Joule heat and
accelerates resistance-mediated degradation.  Aggressive active
cooling consumes additional energy.  A purely safety-focused strategy
derate current so conservatively that deadline requirements are
missed.  Classical constant-current / constant-voltage (CC-CV)
protocols resolve this tension through fixed schedules, ignoring both
real-time cell state and the substantial cell-to-cell variability that
emerges as packs age.

\section{Problem Statement}

Modern BMS research has explored three broad paradigms: rule-based
classical control (CC-CV, multi-stage constant current), Model
Predictive Control (MPC) using Equivalent Circuit Models or
electrochemical models, and single-agent Deep Reinforcement Learning
(DRL) \cite{park2020fastcharging, torchio2016lionsimba, wu2021dqn}.
Each suffers from a fundamental shortcoming.

\begin{enumerate}
  \item \textbf{Classical methods} are static: they cannot adapt to
  individual cell ageing trajectories or dynamic thermal environments.

  \item \textbf{MPC methods} suffer from computational overhead and
  critical sensitivity to model mismatch: as $R_0$ rises with
  degradation, MPC controllers calibrated on fresh cells generate
  systematically incorrect voltage predictions.

  \item \textbf{Single-agent DRL} entangles all objectives into a
  single scalarised reward.  The resulting policy is highly sensitive
  to objective weight selection; small weight perturbations can cause
  policy collapse---catastrophically missing deadlines or violating
  soft voltage boundaries \cite{achiam2017cpo}.
\end{enumerate}

A fourth gap, common to all paradigms, is the absence of
\emph{deterministic} hard safety guarantees: purely neural or
optimisation-based policies have no provable mechanism to
prevent unsafe actions when sensors are faulty or cell parameters
deviate from nominal.

\section{Research Questions}

This thesis addresses four specific research questions derived from
the identified gaps:

\begin{description}
  \item[RQ1.] Does a dedicated health-cap agent, limiting peak current
  to reduce degradation-inducing high-current exposure, produce
  statistically significant capacity-fade reduction compared to an
  independent-deadline baseline under realistic aged-resistance
  variability?

  \item[RQ2.] Can a deterministic analytical safety supervisor
  maintain zero hard voltage and temperature violations across diverse
  sensor fault types (bias, dropout, stuck sensor) at multiple
  charging aggressiveness levels?

  \item[RQ3.] Does a multi-agent functional decomposition
  (with separate proposers for charging speed, thermal regulation, and
  health) outperform single-agent and conventional alternatives across
  a broad range of Dirichlet-sampled objective weights?

  \item[RQ4.] What is the empirically grounded survival distribution
  and Leave-One-Pack-Out (LOPO) RUL prediction accuracy achievable
  from the NASA PCoE ALT dataset, given proper interval-censored
  failure classification?
\end{description}

\section{Contributions}

The principal contributions of this thesis are:

\begin{enumerate}
  \item \textbf{Rigorous empirical dataset audit} of the NASA PCoE
  26-pack ALT dataset, clarifying the CC-only charging protocol
  (3.0\,A, negligible CV taper), correcting the EOL classification
  (9 hard collapses, 5 probable shutoffs, 3 recommissioned, 10
  right-censored), and demonstrating that 67\% of genuine failures
  occur above 0.80 SOH --- invalidating capacity-threshold-only
  prognostic systems.

  \item \textbf{Physics-based simulator calibration} from real
  telemetry: OCV curve, $R_0 = 0.034\,\Omega$ (fresh), thermal
  parameters ($C_\mathrm{th} = 54.5$\,J/K, $h = 0.069$\,W/K,
  $\tau_\mathrm{th} = 791$\,s), and aging-dependent
  $R_0(\mathrm{SOH})$ growth verified across 5{,}479 rest cooling
  blocks.

  \item \textbf{Four-state Adaptive EKF} with state vector
  $[{\rm SOC}, V_1, b, \rho]^\top$, reducing terminal SOC error by
  72.3\% compared to a fixed-parameter EKF and eliminating unsafe
  executed steps on the single-agent tuned policy.

  \item \textbf{MA-BMS architecture} with three cooperative proposer
  agents and a deterministic analytical safety supervisor, achieving
  zero unsafe executed steps under $4\times$ sensor noise and zero
  hard violations across all seven injected sensor fault regimes.

  \item \textbf{Dirichlet weight sensitivity analysis} over
  10{,}000 samples showing that the tuned single-agent policy holds
  an 84.78\% best-weight share, while demonstrating that all four
  methods are Pareto-nondominated, formally qualifying the ranking.

  \item \textbf{Interval-censored survival analysis} and
  Leave-One-Pack-Out RUL study, obtaining Kaplan-Meier median
  throughputs of 630--833\,Ah and Ridge MAE of 122.4\,Ah
  ($R^2 = 0.787$).
\end{enumerate}

\section{Scope and Limitations}

The scope of this work is bounded to the control and prognostics of a
single two-cell in-series pack operating under laboratory accelerated
ageing conditions.  Key limitations are:

\begin{itemize}
  \item Charge-phase current and voltage are not independently logged
  in the NASA dataset; only thermal response during charging is
  validated.  Current dynamics are inferred from the paper (pure CC
  at 3.0\,A).
  \item The RUL sample is small (10 crosser packs in regression);
  LOPO estimates have wide confidence intervals.
  \item The multi-agent coordination layer (coordinator agent) is
  designed but not benchmarked in this thesis; only single-agent and
  independent-baseline variants are fully evaluated.
\end{itemize}

\section{Thesis Organisation}

The remainder of the thesis is organised as follows.
Chapter~\ref{chap.2} reviews the relevant literature.
Chapter~\ref{chap.3} presents the dataset audit.
Chapter~\ref{chap.4} describes simulator calibration and thermal
validation.  Chapter~\ref{chap.5} details the MA-BMS architecture and
EKF design.  Chapter~\ref{chap.6} presents the controller
benchmarking and robustness results.  Chapter~\ref{chap.7} covers
survival analysis and RUL prediction.  Chapter~8 concludes and
outlines future work.
